%%%%%%
%
% $Author: Yogesh$
% $Date: 01.07.2024 $
% $Path: BA24-02-Time-Series\report\Contents\en$
% $Version: 1.0 $
% $Reviewed by: Rakesh $
%
% !TeX encoding = utf8
% !TeX root = Rename
% !TeX TXS-program:bibliography = txs:///bibtex
%
%%%%%%

\chapter{Validation}

\section{General}
Validation is a critical step in ensuring that the XGBoost model for forecasting
air quality in India performs well not only on the training and test datasets but
also in real-world scenarios. This process involves various techniques to
confirm the model's accuracy, reliability, and generalizability over time and
across different conditions \cite{Airqindia:2024}.

\section{Hyperparameter Tuning}
Hyperparameter tuning is an essential part of the validation process. It involves
selecting the optimal set of hyperparameters that maximize the model's performance.
We are now presented with design choices in order to achieve an optimal model
architecture. These choices can be made with the form of parameters, which are
refered to as hyperparameters. Those values are not automatically learned and we
have to tune them. However we don't immediatelly know which parameters to tune
and we may have to explore a huge range of possibilities. So we create a mapping
of hyperparameters and the search space we want to explore.

\bigskip

This is where we can use two popular methods of hyperparameter tuning, the GridSearch
and the RandomSearch. The first option exhaustively searches every possible
combination of those hyperparaters which yields the best result at the cost of
being extremelly slow. The latter option, picks random subsets of the search
space thus being faster and usually providing an adequate result. I'm going to use
RandomSearch for this notebook.For the XGBoost model, key hyperparameters
include learning rate, max depth, subsample, and the number of boosting rounds.
Techniques such as grid search or random search combined with cross-validation
are used to systematically explore and identify the best hyperparameters
\cite{Airqindia:2024}.

\subsection{Grid Search}
Grid search involves exhaustively searching through a predefined set of
hyperparameters. For example, a grid search can be performed over a combination
of learning rates (e.g., 0.01, 0.1), max depths (e.g., 3, 5, 7), and subsample rates
(e.g., 0.8, 1.0). Each combination is evaluated using cross-validation, and the
set that yields the best performance metrics (e.g., lowest RMSE) is selected.

\subsection{Random Search}
Random search is a more efficient alternative to grid search, especially when the
hyperparameter space is large. Instead of evaluating every possible combination,
random search evaluates a fixed number of randomly chosen combinations. This
method can often find a near-optimal solution with less computational effort
\cite{Airqindia:2024}.

\section{Three Ideas for Validation}

\subsection{Streamlit Application for Visualization}
Building a Streamlit application to visualize the predictions and actual air quality
data is an effective way to validate the model. By displaying interactive graphs
and charts, we can observe the model's performance in real-time and identify any
discrepancies between predicted and actual values. This visual approach allows
stakeholders to intuitively understand the model's accuracy and make informed
decisions based on the visualized data.

\subsection{Time Series Split}
Given the temporal nature of air quality data, a time series split validation is
particularly suitable. This technique involves splitting the data sequentially
to mimic real-time prediction scenarios. For example, we can train the model on
data from the first few years and validate it on subsequent years. This approach
ensures that the model can handle temporal dependencies and trends effectively,
providing a realistic evaluation of its forecasting capabilities
\cite{Airqindia:2024}.

\subsection{External Data Validation}
Validating the model using external datasets from different regions or time
periods can further test its generalization capability. By using air quality data
from different locations or new time periods not included in the training set,
we can assess how well the model adapts to new environments and varying
conditions. Successful validation on external datasets would indicate strong robustness
and adaptability, making the model more reliable for practical applications.